\documentclass[11pt]{article}
\usepackage[margin=24mm]{geometry}
\usepackage{amsmath,amssymb,amsthm,hyperref,booktabs}
\newtheorem{theorem}{Theorem}\newtheorem{lemma}{Lemma}
\newcommand{\R}{\mathbb R}\newcommand{\norm}[1]{\lVert#1\rVert}
\newcommand{\ii}{\mathrm i}\newcommand{\perm}{\operatorname{perm}}
\title{A practical local bound at the Fourier matrix}
\author{Research note}\date{7 October 2026}
\begin{document}\maketitle
\noindent\textbf{Scope.} This proves the cutoff-$5$ implication in a specified
neighbourhood. It does not prove that every Hadamard with $M<5$ enters that
neighbourhood. That global confinement is the remaining obligation.
No smoothness or local parametrization of the Hadamard set is assumed.

A raw Hadamard has unit-modulus entries and $HH^*=6I$. Put
$F_{jk}=e^{2\pi\ii jk/6}$, $0\le j,k<6$, and
$M(H)=\max_cR_c(H)$ over all 400 corners. Let $d_F$ be the minimum,
over monomial equivalents of $F$, of the maximum normalized-column
projector distance. A real phase matrix is \emph{horizontal} if its row
and column sums vanish. Norms of phase matrices are Frobenius norms.

\begin{theorem}\label{thm}
Suppose a Hadamard is equivalent to $F\circ e^{\ii u}$ with $u$ horizontal
and $\norm u\le3/10$. Then
\[
 M(H)<5\quad\Longrightarrow\quad
 d_F(H)\le\frac{18471892543}{563200000000}<\frac1{30}.
\]
\end{theorem}

Here is the short proof, with its finite coefficient certificate explained
below. Let $T$ be the horizontal linearized Hadamard kernel at $F$. Its
basis is the horizontal projections of $t_1,t_2,t_1^T,t_2^T$, where $t_1$
is the indicator of odd rows and columns $1,4$, and $t_2$ uses columns
$2,5$. Write $v=P_Tu=T(a,b,c,d)$, and set
\[
 A(x,y)=\max\{|x+y|,|2x-y|,|x-2y|\},\qquad s=A(a,b)+A(c,d).
\]
The exact tangent identity gives $s\le2\norm v\le3/5$.
The following two estimates are the content of the local certificate:
\begin{align}
 \sqrt{M(H)}&\ge 12s-2s^3-29s^4,\label{corner}\\
 d_F(H)&\le \frac{s}{6}
 +\left(\frac1{200}+\frac5{1056}\right)s^3+\frac5{32}s^4.\label{metric}
\end{align}
The polynomial on the right of \eqref{corner} is concave on $s\ge0$.
At both endpoints of $[39/200,3/5]$ it is positive and its square
exceeds $5$ (exact rational comparisons). Thus $M<5$ implies $s<39/200$.
The increasing right side of \eqref{metric}, evaluated at $39/200$, is
exactly the bound in the theorem. This proves the assertion.

\section*{Why the corner estimate is stronger than a generic Taylor bound}
For any corner $B$, the reconstruction supplies a signed real invariant
$d_B$ whose square is the discriminant of the same-row complementary
corner. Therefore $\max_B|d_B|=\sqrt M$.
The first-order corner forms are $0$, the forms $12\ell(a,b)$ and
$12\ell(c,d)$, and the mixed forms $12(\ell(a,b)+\ell'(c,d))$, where
\[
 \ell(x,y)\in\{\pm(x+y),\pm(2x-y),\pm(x-2y)\}.
\]
For each mixed form there are two corners whose linear terms are opposites
and whose effective quadratic terms agree. Their half-difference cancels
the quadratic term. On the sector where its linear term equals $12s$,
the cubic term is at least $-2s^3$. The fourth-order remainder is at most
$29s^4$. Since the absolute value of that half-difference is at most
$\sqrt M$, this proves \eqref{corner}.
This cancellation, rather than a small cutoff, supplies the needed margin.

\section*{Normal displacement: a necessary enclosure for actual Hadamards}
Let $\Phi(u)$ be the thirty real coordinates of the fifteen complex
row overlaps of $F\circ e^{\ii u}$. Let $J=\Phi'(0)$, and $N=T^\perp$
inside the horizontal space. The nonzero eigenvalues of $J^*J$ are
$3,9,12$, with multiplicities $4,12,5$. Hence its inverse from the range
to $N$, extended by zero on the orthogonal complement, has norm
$\norm K=1/\sqrt3$ and $KJ=P_N$.
For horizontal directions, elementary differentiation and Cauchy--Schwarz give
\[
 \norm{\Phi^{(k)}(x)[w_1,\ldots,w_k]}
 \le 2^{(k+1)/2}\prod_j\norm{w_j},\qquad k\ge2.
\]
For $k=2$, this follows from
$\sum_{i<j}\norm{w_i-w_j}^4\le8\norm w^4$; every further derivative
costs at most $\sqrt2\norm{w_j}$. These bounds hold at every base phase.

Solving the quadratic and cubic equations in $N$ gives homogeneous
polynomials $n_2(v),n_3(v)$. With $Q(x,y)=x^2-xy+y^2$, exact calculation gives
\begin{align*}
 \norm v^2&=2(Q(a,b)+Q(c,d)),\\
 \norm{n_2}^2&=\frac43Q(a,b)Q(c,d)\le\frac{s^4}{108},\\
 \norm{n_3}^2&=\frac19Q(a,b)Q(c,d)(Q(a,b)+Q(c,d))\le\frac{s^6}{7776}.
\end{align*}
All Hadamard residuals through degree three vanish at $u_3=v+n_2+n_3$.
For an actual solution, applying $K$ to $\Phi(u)=0$ gives a fixed-point
identity for its normal part. Comparing this identity with $u_3$ yields
\begin{equation}
 \norm{u-u_3}\le\frac38s^4.\label{normal}
\end{equation}
This compares an existing solution with a polynomial; it makes no existence
claim about arbitrary tangent parameters.

For completeness, the rational bounds used to check \eqref{normal} are
as follows. On a closed interval $s\in[l,h]\subset[0,3/5]$, put
\[
 \rho=3/10,\quad c=49/60,\quad c_1=\min(c,\rho/l),\quad
 c_2=1/10,\quad c_3=1/88,\quad t=c_2+c_3h,
\]
with $c_1=c$ if $l=0$. The entire segment joining $u$ and $u_3$ is in radius
\[
 r=\max\{\rho,\min(\rho,ch)+c_2h^2+c_3h^3\}.
\]
The fixed-point derivative is at most $q=2cr<1$. An upper bound for
$\norm{K\Phi(u_3)}/s^4$ is
\[
 B=2cc_1c_3+ct^2+\frac7{18}(3c_1^2t+3c_1t^2h+t^3h^2)
       +\frac{49}{360}(c_1+c_2h+c_3h^2)^4.
\]
Thus the normal error constant is $E=B/(1-q)$. On the adjacent intervals
with endpoints
\[
 0,\quad39/200,\quad8/40,9/40,\ldots,24/40,
\]
the rational checks give $E\le3/8$. At $s=0$ the same contraction
identity directly gives $u=0$.

\section*{The projector-distance estimate}
Let $h_0=(-1,1,-1,1,-1,1)^T/2$, and let $x,w$ be the mean-zero versions
of $(0,a,b,0,a,b)^T$ and $(0,c,d,0,c,d)^T$. Then
$v=h_0x^T+wh_0^T$, with $h_0\perp w$.
A row rephasing centers the rectangle formed by the extreme column vectors.
The resulting column norms divided by $\sqrt6$ are at most $s/6$.
Exact calculation also gives
$h_0^T(n_2)_j=w^T(n_2)_j=0$ for every column $j$.
Consequently the maximum normalized column norm of the balanced $v+n_2$
is at most
\[
 \sqrt{(s/6)^2+\norm{n_2}^2/6}
 \le s/6+s^3/200.
\]
Adding $n_3$ and the error \eqref{normal}, and using $1/\sqrt6<5/12$,
gives \eqref{metric}. This phase bound controls projector distance because
for real column phases $z_j$,
\[
 1-\left|\frac16\sum_j e^{\ii z_j}\right|^2
 =\frac1{36}\sum_{j,k}(1-\cos(z_j-z_k))
 \le\frac16\sum_j(z_j-\bar z)^2.
\]

\section*{The finite corner certificate and its replay}
The normalized sector $s=1$ is one of 36 tetrahedra. Their vertices are
$(v_i,0),(v_{i+1},0),(0,v_j),(0,v_{j+1})$, where the cyclic vertices of
$A(x,y)\le1$ are
\[
 (2,1)/3,(1,2)/3,(-1,1)/3,(-2,-1)/3,(-1,-2)/3,(1,-1)/3.
\]
Cubic Bernstein coefficients on each tetrahedron prove the bound $-2$.
The signed invariant has 48 Laurent terms with coefficients $\pm1$.
Writing it as $\sum_\ell c_\ell\sin(k_\ell\pi/3+\ell u)$ permits a direct
fourth-order remainder bound, term by term.
Bounds on $\ell v$, $\ell n_2$ and $\ell n_3$ come from exact simplex
coefficients; the change from $u_3$ to $u$ uses the normal gradient and
\eqref{normal}. Every acceptance comparison is rational.
The largest resulting corner remainder constant is less than $28.322$,
so the displayed value $29$ is a convenient common upper bound.

The entire certificate, including regeneration of both coefficient tables,
can be replayed with
\begin{verbatim}
python3 -B verify_practical_local.py
\end{verbatim}
The last stage checks the adjacent interval cover and both final rational
inequalities, and writes \texttt{practical\_radius03\_cover.json}.
The input is the defining Hadamard equations and the reconstruction's
complementary-corner identity; MC1, a global classification, and a global
numerical cover are not used in this local theorem.

\section*{The remaining global obligation}
To obtain the requested averaging lemma at cutoff $5$, it remains to prove:
\[
 M(H)<5\quad\Longrightarrow\quad
 H\sim F\circ e^{\ii u}\ \text{with $u$ horizontal and }\norm u\le3/10.
\]
The reconstruction reduces this to excluding strict maximal-corner seeds
outside that phase neighbourhood (with $M=0$ handled by MC1).
A global reconstruction cover or a quantitative critical-strip argument
could establish it. The local theorem above supplies a practical terminal
condition for such a cover; it does not supply the cover itself.
\end{document}
